% ============================================================
\section{Generische Programmierung \& STL (Abschnitt VI)}
% ============================================================
\begin{definition}[: Generische Programmierung]
,,Schreibe Code einmal -- nutze ihn f\"ur viele Typen.'' Algorithmen und Datenstrukturen werden \textbf{unabh\"angig vom konkreten Datentyp} spezifiziert (statt \code{max\_int}, \code{max\_double}, \dots\ nur ein \code{max<T>}).
\end{definition}
Vorteile: \textbf{Wiederverwendung} (DRY), \textbf{Typsicherheit} (Pr\"ufung zur Compile-Zeit), \textbf{Effizienz} (eigene Codeversion je Typ = \emph{Monomorphismus}, kein Laufzeit-Overhead), \textbf{Flexibilit\"at} (Anforderungen an Eigenschaften statt feste Typen).

\subsection{Templates}
Templates sind der zentrale Mechanismus: Code mit Platzhaltern f\"ur Typen/Werte, der \textbf{zur Compile-Zeit} in konkreten Code \"ubersetzt wird.

\subsubsection*{Funktions-Template}
\begin{lstlisting}
template <typename T>
T max(T a, T b) {
    return (a > b) ? a : b;
}
max(3, 5);       // -> max<int>
max(2.5, 1.5);   // -> max<double>
// max(3, 5.5);  // FEHLER: T uneindeutig (int vs double)
\end{lstlisting}
\code{template <...>} steht direkt vor der Funktion und darf nicht abgetrennt werden. Funktioniert nur, wenn die verwendeten Operatoren f\"ur \code{T} existieren.

\subsubsection*{Klassen-Template}
\begin{lstlisting}
template <typename T>
class Box {
    T value;
public:
    Box(T v) : value(v) {}
    T get() { return value; }
};
Box<int> b1(10);   // Typ explizit (noetig vor C++17)
Box b2(3.14);      // ab C++17 deduzierbar
\end{lstlisting}

\subsubsection*{Mehrere Typ- und Nicht-Typ-Parameter}
\begin{lstlisting}
template <typename T, typename U>          // mehrere Typparameter
void printPair(T a, U b) { std::cout << a << ", " << b << '\n'; }

template <typename T, int N>               // Nicht-Typ-Parameter (Wert!)
class Array { public: T data[N]; };
Array<int, 10> a;   // T=int, N=10  (N muss zur Compile-Zeit bekannt sein)
\end{lstlisting}

\subsubsection*{Template-Spezialisierung}
Spezielle Implementierung f\"ur einen konkreten Typ -- hat \textbf{Priorit\"at} vor dem allgemeinen Template.
\begin{lstlisting}
template <typename T> T process(T value) { return value + value; }

template <>                                // Spezialisierung fuer bool
bool process<bool>(bool value) { return !value; }
\end{lstlisting}

\begin{merke}
Templates werden erst \textbf{beim Verwenden} instanziiert und f\"ur jede \"Ubersetzungseinheit separat erzeugt. Deshalb m\"ussen Templates \textbf{inkl. ihrer Logik in der Header-Datei (\code{.hpp})} stehen -- Implementierung in einer \code{.cpp} f\"uhrt zu \emph{Undefined Behaviour}.
\end{merke}

\subsection{STL -- Standard Template Library}
Die STL kombiniert \textbf{Container} (Objektsammlungen), \textbf{Iteratoren} (einheitlicher Durchlauf) und \textbf{Algorithmen} (Such-/Sortierfunktionen). Vorteile: optimierte Performance, weniger Speicherfehler, standardisierte Interfaces.

\subsubsection*{Container-Kategorien}
\begin{center}\small
\begin{tabular}{@{}ll@{}}
\toprule
\textbf{Kategorie} & \textbf{Beschreibung / Beispiele} \\
\midrule
Sequence & geordnet, Zugriff per Index/Iterator (\code{vector}, \code{list}, \code{array}) \\
Adapter & Wrapper, der einen Container einschr\"ankt (\code{stack}, \code{queue}) \\
Associative & automatisch sortiert, Zugriff per Key, Baum (\code{map}, \code{set}) \\
Unordered & Key-basiert, unsortiert, Hash (\code{unordered\_map}) \\
\bottomrule
\end{tabular}
\end{center}

\subsubsection*{Die wichtigsten Container}
\begin{center}\small
\begin{tabular}{@{}lll@{}}
\toprule
 & \textbf{\code{std::vector}} & \textbf{\code{std::map}} \\
\midrule
Modell & dynamisches Array & Key$\to$Value (sortiert) \\
Speicher & zusammenh\"angend & Baumstruktur \\
Indexzugriff & \code{v[i]} -- $O(1)$ & --- \\
Suche & $O(n)$ & per Key $O(\log n)$ \\
typisch & \code{push\_back}, \code{front}, \code{back} & \code{m["k"]=v}, \code{insert}, \code{erase} \\
\bottomrule
\end{tabular}
\end{center}
\code{std::list} ist eine \textbf{doppelt verkettete Liste}: kein Indexzugriff, aber schnelles Einf\"ugen/L\"oschen an beliebiger Stelle (\code{push\_front}, \code{push\_back}).
\begin{lstlisting}
std::vector<int> v = {1, 2, 3};
v.push_back(4);                // anhaengen
std::cout << v[1] << v.front() << v.back();

std::map<std::string, int> m;
m["Alice"] = 25;               // anlegen oder aendern
m.insert({"Bob", 30});
m.erase("Alice");
\end{lstlisting}
Gemeinsame Konzepte aller Container: \code{c.size()}, \code{c.empty()}, \code{c.insert(...)}, \code{c.erase(...)}, for-each-Schleife und STL-Algorithmen.

\subsection{for-each-Schleife}
\begin{lstlisting}
std::vector<int> v = {1, 2, 3};
for (int x : v)        { std::cout << x << " "; }   // Kopie
for (auto& x : v)      { x *= 2; }                  // Referenz: veraendert v
\end{lstlisting}
Durchl\"auft alle Elemente; kein Indexzugriff, Reihenfolge/Richtung nicht anpassbar.

\subsection{Iteratoren}
Ein \emph{Iterator} ist ein \textbf{verallgemeinerter Zeiger}: zeigt auf ein Element, kann weiterbewegt werden und erlaubt Zugriff. Die for-each-Schleife nutzt intern Iteratoren.
\begin{lstlisting}
for (std::vector<int>::iterator it = v.begin(); it != v.end(); ++it) {
    std::cout << *it << " ";    // *it: Zugriff,  ++it: naechstes Element
}
\end{lstlisting}
\code{begin()} liefert den ersten, \code{end()} zeigt \textbf{hinter} das letzte Element -- der Bereich ist halboffen \code{[begin, end)}.

\subsubsection*{Iterator-Kategorien (hierarchisch, je au\ss en st\"arker)}
\begin{center}\small
\begin{tabular}{@{}llll@{}}
\toprule
\textbf{Kategorie} & \textbf{Kann} & \textbf{Operationen} & \textbf{Container} \\
\midrule
Input / Output & einmal lesen \emph{bzw.}\ schreiben & \code{*i}, \code{++} & Streams \\
Forward & mehrfach lesen+schreiben & \code{*i}, \code{++} & \code{forward\_list} \\
Bidirectional & vorw\"arts \& r\"uckw\"arts & \code{++}, \code{--} & \code{list}, \code{map} \\
Random-Access & beliebiger Index & \code{++ -- += []} \code{< >} & \code{vector} \\
\bottomrule
\end{tabular}
\end{center}
\begin{merke}
Beim Iterator: \code{.} greift auf Member des \emph{Iterators} zu, \code{->} auf Member des \emph{Elements}. (\code{it->name} statt \code{(*it).name}.)
\end{merke}

\subsection{STL-Algorithmen}
Generische Funktionen, die auf \textbf{Bereichen} \code{[begin, end)} arbeiten -- \emph{nicht} auf Containern. Dadurch containerunabh\"angig.
\begin{syntaxbox}
\code{std::algorithmus(begin, end, ...);}
\end{syntaxbox}
\begin{center}\small
\begin{tabular}{@{}ll@{}}
\toprule
\textbf{Algorithmus} & \textbf{Wirkung} \\
\midrule
\code{std::sort} & sortiert aufsteigend ($O(n\log n)$) \\
\code{std::find} & erstes passendes Element \\
\code{std::count} & Anzahl eines Werts \\
\code{std::accumulate} & Summe (aus \code{<numeric>}) \\
\code{std::copy} / \code{transform} & kopieren / Funktion anwenden \\
\code{std::reverse} / \code{replace} & umkehren / Werte ersetzen \\
\bottomrule
\end{tabular}
\end{center}
\begin{lstlisting}
#include <algorithm>
std::vector<int> v = {4, 1, 7, 3};
std::sort(v.begin(), v.end());                 // 1 3 4 7
auto it = std::find(v.begin(), v.end(), 3);
if (it != v.end()) std::cout << "gefunden: " << *it;
\end{lstlisting}

\begin{tipp}
\textbf{Trade-offs:} Templates erzeugen oft sehr lange, schwer lesbare Compiler-Fehler, h\"ohere Codekomplexit\"at sowie l\"angere Compile-Zeiten und gr\"o\ss ere Bin\"ardateien (\emph{Code-Bloat}), da pro Typ eigener Code entsteht.
\end{tipp}
\clearpage
